\chapter{Symplectic Geometry}
\label{chap:SymplecticGeometry}
\section{Hamiltonian Mechanics}
\subsubsection*{Hamilton's Equations}
Hamiltonian mechanics is a reformulation of classical mechanics. 
It arises from Lagrangian mechanics by a Legendre transformation. 
The state of a mechanical system is described by a point in a 
$2n$-dimensional phase space, with coordinates $(q_1,\ldots,q_n,
p_1,\ldots,p_n)$, where the $q_i$ are the generalized coordinates,
and the $p_i$ are the conjugate momenta. The time evolution of the
system is given by \textbf{Hamilton's equations}:
\begin{equation*}
    \left\{
    \begin{aligned}
        \pd{q_i}{t} &= \pd{H}{p_i}, \\
        \pd{p_i}{t} &= -\pd{H}{q_i},
    \end{aligned}
    \quad i=1,\ldots,n,\right.
\end{equation*}
where $H(q,p,t)$ is the \textbf{Hamiltonian function}, which is usually the
total energy of the system.

Suppoese $(q(t),p(t))$ is a solution to Hamilton's equations. 
Then for any function $f(q,p,t)$, its time evalution is given by 
\begin{equation*}
    \frac{\md f}{\md t} = \pd{f}{t} + \sum_{i=1}^n\left(
    \pd{f}{q_i}\frac{\md q_i}{\md t} + 
    \pd{f}{p_i}\frac{\md p_i}{\md t}\right)
    = \pd{f}{t} + \sum_{i=1}^{n}\left(\pd{f}{q_i}\pd{H}{p_i} 
    - \pd{f}{p_i}\pd{H}{q_i}\right).
\end{equation*}
We define the \textbf{Poisson bracket} of two functions $f$ and $g$ by
\begin{equation*}
    \{f,g\} = \sum_{i=1}^{n}\left(\pd{f}{q_i}\pd{g}{p_i} 
    - \pd{f}{p_i}\pd{g}{q_i}\right).
\end{equation*}
Then the time evolution of $f$ can be written as 
\begin{equation*}
    \frac{\md f}{\md t} = \pd{f}{t} + \{f,H\}.
\end{equation*}
So $f$ is a constant of motion if and only if $\pd{f}{t} + \{f,H\} = 0$. 
If in particular $f$ does not depend on $t$, then $f$ is a constant of 
motion if and only if $\{f,H\} = 0$. 

\subsubsection*{Hamilton's Equations and Symplectic Forms}
To formulate Hamiltonian mechanics in a symplectic geometric framework,
we introduce the symplectic form $\omega$ on the phase space:
\begin{equation*}
    \omega = \sum_{i=1}^n \md q_i \wedge \md p_i.
\end{equation*}
This is a closed, non-degenerate 2-form. Let $H = H(q,p)$ be a 
Hamiltonian function. The Hamiltonian vector field $X_H$ associated 
to $H$ is defined by the equation
\begin{equation*}
    \iota_{X_H}\omega = \md H, 
\end{equation*}
where $\iota_{X_H}$ denotes the interior product with $X_H$. 
The left-hand side is 
\begin{equation*}
    \iota_{X_H}\omega = \sum_{i=1}^n\Big(\md q_i(X_H)\md p_i 
    - \md p_i(X_H)\md q_i\Big).
\end{equation*}
The right-hand side is
\begin{equation*}
    \md H = \sum_{i=1}^n\left(\pd{H}{q_i}\md q_i + \pd{H}{p_i}\md p_i\right).
\end{equation*}
So we have
\begin{equation*}
    \md q_i(X_H) = \pd{H}{p_i}, \quad \md p_i(X_H) = -\pd{H}{q_i}.
\end{equation*}
Thus the Hamiltonian vector field is 
\begin{equation*}
    X_H = \sum_{i=1}^n\left(\pd{H}{p_i}\pd{}{q_i} 
    - \pd{H}{q_i}\pd{}{p_i}\right).
\end{equation*}
The time evolution of the system is given by the flow of $X_H$, which 
is exactly Hamilton's equations. 

For any smooth functions $f$, we can uniquely determine a verctor field 
$X_f$ by solving the equation 
\begin{equation*}
    \iota_{X_f}\omega = \md f.
\end{equation*}
For any two smooth functions $f,g$, We can define the Poission bracket to be 
\begin{equation*}
    \{f,g\} := \omega(X_f,X_g).
\end{equation*}

When $\omega = \sum_{i=1}^{n}\md q_i\wedge \md p_i$, we can get
\begin{equation*}
    X_f = \sum_{i=1}^{n}\left(\pd{f}{p_i}\pd{}{q_i} 
    - \pd{f}{q_i}\pd{}{p_i}\right).
\end{equation*}
So 
\begin{align*}
    \omega(X_f,X_g) &= \sum_{i,j=1}^{n}
    \omega\left(\pd{f}{p_i}\pd{}{q_i}-\pd{f}{q_i}\pd{}{p_i},
    \pd{g}{p_j}\pd{}{q_j}-\pd{g}{q_j}\pd{}{p_j}\right) \\
    &= \sum_{i,j=1}^{n}\left(\pd{f}{q_i}\pd{g}{p_j} 
    - \pd{f}{p_i}\pd{g}{q_j}\right)\delta_{ij} \\
    &= \sum_{i=1}^{n}
    \left(\pd{f}{q_i}\pd{g}{p_i} - \pd{f}{p_i}\pd{g}{q_i}\right).
\end{align*}

\subsubsection*{Canonical Transformations}
Let $(Q(q,p),P(q,p))$ be a transformation from $(q,p)$ to $(Q,P)$. We say
it is a \textbf{canonical transformation} if it preserves the form of 
Hamilton's equations, i.e., if $(Q(t),P(t))$ satisfies
\begin{equation*}
    \left\{
    \begin{aligned}
        \pd{Q_i}{t} &= \pd{K}{P_i}, \\
        \pd{P_i}{t} &= -\pd{K}{Q_i},
    \end{aligned}
    \quad i=1,\ldots,n,\right.
\end{equation*}
for some Hamiltonian function $K(Q,P,t)$. This is equivalent to the 
\textbf{Poisson bracket condition}:
\begin{gather*}
    \{P_i,Q_j\}_{q,p} = \delta_{ij}, \\
    \{Q_i,Q_j\}_{q,p} = 0, \\
    \{P_i,P_j\}_{q,p} = 0.
\end{gather*}
Another equivalent condition is the \textbf{differential form condition}:
\begin{equation*}
    \sum_{i=1}^n\md Q_i\wedge\md P_i=\sum_{i=1}^n\md q_i\wedge\md p_i.
\end{equation*}

\subsubsection*{Canonical transformations and Symplectic Forms}
Let $(Q(q,p),P(q,p))$ be a transformation from $(q,p)$ to $(Q,P)$. The 
Jacobian matrix of this transformation is
\begin{equation*}
    J = \frac{\p(Q,P)}{\p(q,p)} = 
    \begin{pmatrix}
        \pd{Q}{q} & \pd{Q}{p} \\ \pd{P}{q} & \pd{P}{p}
    \end{pmatrix} =: 
    \begin{pmatrix}
        A & B \\ C & D
    \end{pmatrix}.
\end{equation*}
\begin{theorem}
    The Poission bracket condition is equivalent to $J^T$ being symplectic.
\end{theorem}
\begin{proof}
    The Poission bracket condition is
    \begin{gather*}
        \{P_k,Q_l\}_{q,p} = \sum_{i=1}^{n}\Big(\pd{Q_l}{q_i}\pd{P_k}{p_i}
        -\pd{Q_l}{p_i}\pd{P_k}{q_i}\Big) = \delta_{kl}, \\
        \{Q_k,Q_l\}_{q,p} = \sum_{i=1}^{n}\Big(\pd{Q_k}{q_i}\pd{Q_l}{p_i}
        -\pd{Q_k}{p_i}\pd{Q_l}{q_i}\Big) = 0, \\
        \{P_k,P_l\}_{q,p} = \sum_{i=1}^{n}\Big(\pd{P_k}{q_i}\pd{P_l}{p_i}
        -\pd{P_k}{p_i}\pd{P_l}{q_i}\Big) = 0. 
    \end{gather*}
    Written in a tigher form, they are
    \begin{gather*}
        \pd{Q}{q}\Big(\pd{P}{p}\Big)^T-\pd{Q}{p}\Big(\pd{P}{q}\Big)^T = I 
        \Rightarrow A D^T - B C^T = I, \\
        \pd{Q}{q}\Big(\pd{Q}{p}\Big)^T-\pd{Q}{p}\Big(\pd{Q}{q}\Big)^T = 0 
        \Rightarrow A B^T - B A^T = 0, \\
        \pd{P}{q}\Big(\pd{P}{p}\Big)^T-\pd{P}{p}\Big(\pd{P}{q}\Big)^T = 0 
        \Rightarrow C D^T - D C^T = 0.
        \end{gather*}
    This is equivalent to say 
    \begin{align*}
        J \Omega J^T &= 
        \begin{pmatrix}
            A & B \\ C & D
        \end{pmatrix}
        \begin{pmatrix}
            0 & I \\ -I & 0
        \end{pmatrix}
        \begin{pmatrix}
            A^T & C^T \\ B^T & D^T
        \end{pmatrix}  \\ 
        &= 
        \begin{pmatrix}
            A B^T - B A^T & A D^T - B C^T \\
            C B^T - D A^T & C D^T - D C^T
        \end{pmatrix} \\
        &= 
        \begin{pmatrix}
            0 & I \\ -I & 0
        \end{pmatrix} = \Omega.
    \end{align*}
\end{proof}
\begin{theorem}
    The differential form condition is equivalent to $J$ being symplectic.
\end{theorem}
\begin{proof}
    Let $\omega_{(q,p)} = \sum\limits_{i=1}^n \md q_i\wedge \md p_i$,
    $\omega_{(Q,P)} = \sum\limits_{i=1}^n \md Q_i\wedge \md P_i$. Then 
    \begin{align*}
        \omega_{(Q,P)} &= 
        \sum_{i=1}^n\left(\sum_{k=1}^{n}\pd{Q_i}{q_k}\md q_k 
        + \pd{Q_i}{p_k}\md p_k\right) \wedge 
        \left(\sum_{l=1}^{n}\pd{P_i}{q_l}\md q_l 
        + \pd{P_i}{p_l}\md p_l\right) \\ 
        &= \sum_{i=1}^n\sum_{k,l=1}^{n}\left(\pd{Q_i}{q_k}\pd{P_i}{p_l}
        -\pd{Q_i}{p_l}\pd{P_i}{q_k}\right)\md q_k\wedge \md p_l \\
        &\quad + \sum_{i=1}^n\sum_{k<l}\left(\pd{Q_i}{q_k}\pd{P_i}{q_l}
        -\pd{Q_i}{q_l}\pd{P_i}{q_k}\right)\md q_k\wedge \md q_l \\
        &\quad + \sum_{i=1}^n\sum_{k<l}\left(\pd{Q_i}{p_k}\pd{P_i}{p_l}
        -\pd{Q_i}{p_l}\pd{P_i}{p_k}\right)\md p_k\wedge \md p_l.
    \end{align*}
    Thus $\omega_{(Q,P)} = \omega_{(q,p)}$ if and only if 
    \begin{gather*}
        \sum_{i=1}^n\left(\pd{Q_i}{q_k}\pd{P_i}{p_l}
        -\pd{Q_i}{p_l}\pd{P_i}{q_k}\right) = \delta_{kl}, \\
        \sum_{i=1}^n\left(\pd{Q_i}{q_k}\pd{P_i}{q_l}
        -\pd{Q_i}{q_l}\pd{P_i}{q_k}\right) = 0, \\
        \sum_{i=1}^n\left(\pd{Q_i}{p_k}\pd{P_i}{p_l}
        -\pd{Q_i}{p_l}\pd{P_i}{p_k}\right) = 0. 
    \end{gather*}
    Written in a tigher form, they are 
    \begin{gather*}
        \Big(\pd{Q}{q}\Big)^T\pd{P}{p}-\Big(\pd{P}{q}\Big)^T\pd{Q}{p} = I
        \Rightarrow A^{T}D - C^{T}B = I, \\
        \Big(\pd{Q}{q}\Big)^T\pd{Q}{p}-\Big(\pd{P}{q}\Big)^T\pd{P}{p} = 0
        \Rightarrow A^{T}B - C^{T}A = 0, \\
        \Big(\pd{Q}{p}\Big)^T\pd{P}{q}-\Big(\pd{P}{p}\Big)^T\pd{Q}{q} = 0
        \Rightarrow B^{T}D - D^{T}B = 0.
    \end{gather*}
    This is equivalent to say 
    \begin{align*}
        J^T \Omega J &= 
        \begin{pmatrix}
            A^T & C^T \\ B^T & D^T
        \end{pmatrix}
        \begin{pmatrix}
            0 & I \\ -I & 0
        \end{pmatrix}
        \begin{pmatrix}
            A & B \\ C & D
        \end{pmatrix}  \\ 
        &= 
        \begin{pmatrix}
            A^{T}C - C^{T}A & A^{T}D - C^{T}B \\
            B^{T}C - D^{T}A & B^{T}D - D^{T}B
        \end{pmatrix} \\
        &= 
        \begin{pmatrix}
            0 & I \\ -I & 0
        \end{pmatrix} = \Omega.
    \end{align*}
\end{proof}
\begin{theorem}
    $J$ is symplectic if and only if $J^T$ is symplectic.
\end{theorem}
\begin{proof}
    \begin{align*}
        J \text{ is symplectic} &\Leftrightarrow J^T\Omega J = \Omega \\ 
        &\Leftrightarrow J^{-1}\Omega (J^T)^{-1} = \Omega \\ 
        &\Leftrightarrow J\Omega J^T = \Omega \\
        &\Leftrightarrow J^T \text{ is symplectic}.
    \end{align*} 
\end{proof}
Those three theorems together give the following conclusion:
\begin{theorem}
    The Poission bracket condition is equivalent to the differential form 
    condition.
\end{theorem}